\textbf{Phần 1. Trắc nghiệm lựa chọn}
\setcounterpageref{ex}{0}

\begin{ex} %[3D13N68]
Đơn vị tiền tệ của nước Việt Nam là
\choice
{Đô la}
{Xu}
{\True Đồng}
{Hào}
\loigiai{
Đơn vị tiền tệ là Đồng.\\
\textbf{Đáp án của câu hỏi là C}
}
\end{ex}

\begin{ex} %[3D13N68]
Tờ tiền nào sau đây có mệnh giá lớn nhất?
\choice
{$10,000$ đồng}
{$50,000$ đồng}
{$20,000$ đồng}
{\True $100,000$ đồng}
\loigiai{
$100,000$ đồng có mệnh giá lớn nhất.\\
\textbf{Đáp án của câu hỏi là D}
}
\end{ex}

\begin{ex} %[3D13H68]
Hai tờ tiền $50,000$ đồng có tổng giá trị bằng
\choice
{\True $100,000$ đồng}
{$70,000$ đồng}
{$80,000$ đồng}
{$90,000$ đồng}
\loigiai{
$50,000 \times 2 = 100,000$ đồng.\\
\textbf{Đáp án của câu hỏi là A}
}
\end{ex}

\begin{ex} %[3D13H68]
Năm tờ tiền $10,000$ đồng đổi được một tờ tiền mệnh giá
\choice
{$20,000$ đồng}
{$100,000$ đồng}
{\True $50,000$ đồng}
{$5,000$ đồng}
\loigiai{
$10,000 \times 5 = 50,000$ đồng.\\
\textbf{Đáp án của câu hỏi là C}
}
\end{ex}

\begin{ex} %[3D13H68]
Bạn Nam có 3 tờ $20,000$ đồng và 1 tờ $10,000$ đồng. Nam có tất cả số tiền là
\choice
{\True $70,000$ đồng}
{$60,000$ đồng}
{$80,000$ đồng}
{$50,000$ đồng}
\loigiai{
$20,000 \times 3 + 10,000 = 70,000$ đồng.\\
\textbf{Đáp án của câu hỏi là A}
}
\end{ex}

\begin{ex} %[3D13H68]
Một cây bút chì giá $5,000$ đồng, một quyển vở giá $10,000$ đồng. Mua cả hai thứ hết số tiền là
\choice
{$12,000$ đồng}
{\True $15,000$ đồng}
{$20,000$ đồng}
{$18,000$ đồng}
\loigiai{
$5,000 + 10,000 = 15,000$ đồng.\\
\textbf{Đáp án của câu hỏi là B}
}
\end{ex}

\begin{ex} %[3D13V68]
Mai có $50,000$ đồng, Mai mua một chiếc thước kẻ giá $8,000$ đồng và một hộp bút màu giá $25,000$ đồng. Mai còn lại số tiền là
\choice
{$16,000$ đồng}
{$18,000$ đồng}
{$20,000$ đồng}
{\True $17,000$ đồng}
\loigiai{
Đã mua: $8,000 + 25,000 = 33,000$ đồng.\\
Còn lại: $50,000 - 33,000 = 17,000$ đồng.\\
\textbf{Đáp án của câu hỏi là D}
}
\end{ex}

\begin{ex} %[3D13V68]
Mẹ đưa cho cô bán hàng tờ $100,000$ đồng để mua một hộp sữa giá $65,000$ đồng. Cô bán hàng phải trả lại mẹ số tiền là
\choice
{$30,000$ đồng}
{\True $35,000$ đồng}
{$40,000$ đồng}
{$45,000$ đồng}
\loigiai{
$100,000 - 65,000 = 35,000$ đồng.\\
\textbf{Đáp án của câu hỏi là B}
}
\end{ex}

\begin{ex} %[3D13V68]
Mua 4 cái bánh, mỗi cái giá $12,000$ đồng. Tổng số tiền phải trả là
\choice
{$44,000$ đồng}
{$50,000$ đồng}
{\True $48,000$ đồng}
{$52,000$ đồng}
\loigiai{
$12,000 \times 4 = 48,000$ đồng.\\
\textbf{Đáp án của câu hỏi là C}
}
\end{ex}

\begin{ex} %[3D13V68]
Có $60,000$ đồng mua được bao nhiêu quyển vở loại $15,000$ đồng một quyển?
\choice
{$3$ quyển}
{\True $4$ quyển}
{$5$ quyển}
{$6$ quyển}
\loigiai{
$60,000 : 15,000 = 4$ quyển.\\
\textbf{Đáp án của câu hỏi là B}
}
\end{ex}

\begin{ex} %[3D13V68]
Một gói bánh giá $24,000$ đồng, đắt gấp đôi một gói kẹo. Một gói kẹo có giá là
\choice
{$10,000$ đồng}
{$14,000$ đồng}
{$16,000$ đồng}
{\True $12,000$ đồng}
\loigiai{
$24,000 : 2 = 12,000$ đồng.\\
\textbf{Đáp án của câu hỏi là D}
}
\end{ex}

\begin{ex} %[3D13C68]
Hùng mua 3 cây bút mực và đưa tờ tiền $50,000$ đồng, được thối lại $14,000$ đồng. Giá tiền của mỗi cây bút mực là
\choice
{\True $12,000$ đồng}
{$11,000$ đồng}
{$13,000$ đồng}
{$14,000$ đồng}
\loigiai{
3 cây bút giá: $50,000 - 14,000 = 36,000$ đồng.\\
Mỗi cây bút giá: $36,000 : 3 = 12,000$ đồng.\\
\textbf{Đáp án của câu hỏi là A}
}
\end{ex}

\textbf{Phần 2. Trắc nghiệm đúng sai}
\setcounterpageref{ex}{0}

\begin{ex} %[3D13H68]
Xét tính đúng sai về tiền tệ:
\choiceTF
{\True 2 tờ $50,000$ đồng đổi được 1 tờ $100,000$ đồng.}
{\True 5 tờ $20,000$ đồng có tổng giá trị $100,000$ đồng.}
{1 tờ $100,000$ đồng đổi được 4 tờ $20,000$ đồng.}
{\True 1 tờ $100,000$ đồng đổi được 5 tờ $20,000$ đồng.}
\loigiai{
\begin{itemize}
    \item $50000 \times 2 = 100000$: Đúng.
    \item $20000 \times 5 = 100000$: Đúng.
    \item $100000 : 20000 = 5$ tờ (không phải 4 tờ). Ý 3 Sai, Ý 4 Đúng.
\end{itemize}
\begin{itemize}
    \item a) \textbf{Đúng}.
    \item b) \textbf{Đúng}.
    \item c) \textbf{Sai}.
    \item d) \textbf{Đúng}.
\end{itemize}
}
\end{ex}

\begin{ex} %[3D13V68]
Bạn Lan có $80,000$ đồng. Lan mua 1 cuốn truyện giá $45,000$ đồng và 2 cây bút chì giá $10,000$ đồng một cây. Xét tính đúng sai:
\choiceTF
{\True Tiền mua 2 cây bút chì là $20,000$ đồng.}
{\True Tổng số tiền Lan đã mua sắm là $65,000$ đồng.}
{Lan còn lại $20,000$ đồng.}
{\True Lan còn lại $15,000$ đồng ($80000 - 65000 = 15000$).}
\loigiai{
\begin{itemize}
    \item Bút: $10000 \times 2 = 20000$ đồng. Đúng.
    \item Tổng: $45000 + 20000 = 65000$ đồng. Đúng.
    \item Còn lại: $80000 - 65000 = 15000$ đồng $\neq 20000$. Ý 3 Sai, Ý 4 Đúng.
\end{itemize}
\begin{itemize}
    \item a) \textbf{Đúng}.
    \item b) \textbf{Đúng}.
    \item c) \textbf{Sai}.
    \item d) \textbf{Đúng}.
\end{itemize}
}
\end{ex}

\textbf{Phần 3. Trả lời ngắn}
\setcounterpageref{ex}{0}

\begin{ex} %[3D13H68]
Tính tổng số tiền gồm 4 tờ $20,000$ đồng và 1 tờ $10,000$ đồng theo nghìn đồng (ví dụ 90000 đồng ghi 90000).
\shortans{90000}
\loigiai{
\textbf{Đáp số:} $90000$\\
$20000 \times 4 + 10000 = 90000$ đồng.
}
\end{ex}

\begin{ex} %[3D13H68]
Cần bao nhiêu tờ tiền $5,000$ đồng để đổi được một tờ $20,000$ đồng?
\shortans{4}
\loigiai{
\textbf{Đáp số:} $4$\\
$20000 : 5000 = 4$ tờ.
}
\end{ex}

\begin{ex} %[3D13V68]
Mua 3 quyển truyện tranh, mỗi quyển giá $18,000$ đồng. Tính tổng số tiền phải trả.
\shortans{54000}
\loigiai{
\textbf{Đáp số:} $54000$\\
$18000 \times 3 = 54000$ đồng.
}
\end{ex}

\begin{ex} %[3D13V68]
Bác An đưa tờ $100,000$ đồng mua xăng hết $72,000$ đồng. Người bán hàng phải thối lại bao nhiêu đồng?
\shortans{28000}
\loigiai{
\textbf{Đáp số:} $28000$\\
$100000 - 72000 = 28000$ đồng.
}
\end{ex}

\begin{ex} %[3D13V68]
Một hộp bút màu giá $35,000$ đồng, một cái cặp sách giá $60,000$ đồng. Hỏi mua cả hai thứ hết bao nhiêu đồng?
\shortans{95000}
\loigiai{
\textbf{Đáp số:} $95000$\\
$35000 + 60000 = 95000$ đồng.
}
\end{ex}

\begin{ex} %[3D13C68]
Một bạn học sinh mua 5 quyển vở cùng loại và đưa tờ $100,000$ đồng, được trả lại $35,000$ đồng. Hỏi mỗi quyển vở giá bao nhiêu đồng?
\shortans{13000}
\loigiai{
\textbf{Đáp số:} $13000$\\
5 quyển vở giá: $100000 - 35000 = 65000$ đồng.\\
Mỗi quyển giá: $65000 : 5 = 13000$ đồng.
}
\end{ex}
